\section{«Малая» и «большая» конечные арифметики. Примеры.}

Ф.А. Новиков выделяет два уровня формализации вычислений в ЭВМ, называя их «малой» и «большой» арифметиками.

\subsection*{«Малая» конечная арифметика}
Это арифметика на уровне \textbf{отдельных битов}.
\begin{itemize}
    \item \textbf{Носитель:} Булево множество $\mathbb{B} = \{0, 1\}$.
    \item \textbf{Операции:} Логическое И ($\land$), ИЛИ ($\lor$), НЕ ($\neg$), Исключающее ИЛИ ($\oplus$).
    \item \textbf{Суть:} Описывает работу логических вентилей процессора. Здесь нет понятия «числа» в привычном смысле, только логические значения.
\end{itemize}

\subsection*{«Большая» конечная арифметика}
Это арифметика на уровне \textbf{машинных слов} (регистров).
\begin{itemize}
    \item \textbf{Носитель:} Множество всех возможных значений регистра $M = \{0, 1, \dots, 2^k-1\}$.
    \item \textbf{Операции:} Сложение и умножение по модулю $2^k$.
    \item \textbf{Суть:} Описывает поведение целых чисел в программах. Это кольцо $(\mathbb{Z}_{2^k}; +, \ast)$. 
\end{itemize}

\subsection*{Связь между ними}
«Большая» арифметика реализуется через «малую». Например, операция сложения в «большой» арифметике строится из логических операций «малой» арифметики (сумматоров), учитывающих \textbf{перенос} (carry).

\paragraph{Примеры}
\begin{enumerate}
    \item \textbf{Малая:} Вычисление четности бита через $a \oplus b$.
    \item \textbf{Большая:} Работа счетчика цикла, который при достижении значения $2^k-1$ обнуляется при следующем шаге (свойство кольца вычетов).
    \item \textbf{Логический сдвиг:} Операция в большой арифметике, которая эффективно соответствует умножению на 2 в кольце, но реализуется как перемещение битов.
\end{enumerate}